\scanpage{155}
\unsure{ramène|revenir} au cas où $G$ \unsure{agit|opère} transitivement sur l'ens des comp. connexes. Soit $X_0$ \inserted{($0\in I$)} une d'elles, $G_0$ son stabilisateur, alors il est évident que $X$ s'identifie à $X\times_{G_0}G$, et $X/G$ à $X_0/G_0$. \deleted{Enfin, alors} on est ramené à $X_0$ connexe.

\deleted{Soit $H$ \unsure{$x\in X_0$}, soit $\unsure{G_{xi}|G_{x_i}|G_i}$ son groupe d'inertie, il opère trivial\textsuperscript{t} sur $X$, \unsure{alors est sous} l'ens}

On \unsure{peut|va} supposer que $G_0$ opère \uline{fidèlt}. Mais \uline{alors il opère sans inertie} donc $X\longrightarrow X/G$ est \uline{étale} (prop \unsure{--}). Reste à voir que $X/G$ est étale sur $Y$.

Faisant une extension fid plate de la base $Y$ convenable, on est ramené au cas où l'ext\textsuperscript{n} résiduelle $k(x)/k(y)$ est triviale. \unsure{Concluons|Conclusion} on voit sur le diagramme $\widehat{\underline{\mathcal O}}_y\xrightarrow{\unsure{-|\sim}}\widehat{\underline{\mathcal O}}_{x'}$, cf le

\diagram{\begin{tikzcd}[ampersand replacement=\&,row sep=small,column sep=small]
x\&x'\\
X\arrow[dr]\&\\
\&X'=X/G\arrow[d]\\
\&Y\ni y
\end{tikzcd}}

\diagram{\begin{tikzcd}[ampersand replacement=\&,row sep=small,column sep=tiny]
k(x')\arrow[dr,dash]\&\\
\&k(x)\arrow[d,dash]\\
\&k(y)
\end{tikzcd}}

\diagram{\begin{tikzcd}[ampersand replacement=\&,column sep=small]
\widehat{\underline{\mathcal O}}_y\arrow[r]\arrow[rr,bend right=18,"\sim"']\&\widehat{\underline{\mathcal O}}_{x'}\arrow[r,"\sim"]\&\widehat{\underline{\mathcal O}}_x
\end{tikzcd}}

\uline{Corollaire}

Soit $X$ étale sur $Y$ en $x$, $g:X\longrightarrow X'$ un $Y$-morphisme étale en $x$, alors $X'$ est étale sur $Y$ en $g(x)$.

% archive sheet number: 154

